% Part: lambda-calculus
% Chapter: introduction
% Section: computable-lambda

\documentclass[../../../include/open-logic-section]{subfiles}

\begin{document}

\olfileid{lam}{int}{lrp}
\olsection{સંગણનીય વિધેયો \usetoken{S}{lambda definable} છે}

\begin{thm}
\ollabel{thm:computable-lambda}
દરેક સંગણનીય આંશિક વિધેય !!{lambda definable} છે.
\end{thm}

\begin{proof}
દરેક આંશિક સંગણનીય વિધેય~$f$ કોઈ લૅમ્બડા પદ~$F$ વડે
!!{lambda defined} છે એમ બતાવવું છે. ક્લીનીના સામાન્ય સ્વરૂપના
પ્રમેય મુજબ, દરેક આદિમ પુનરાવર્તી વિધેય કોઈ લૅમ્બડા પદ વડે
!!{lambda defined} છે અને !!{lambda definable} વિધેયો યોગ્ય
સંયોજનો તથા અમર્યાદિત શોધ હેઠળ બંધ છે એમ બતાવવું પૂરતું છે.
દરેક આદિમ પુનરાવર્તી વિધેય લૅમ્બડા પદ વડે !!{lambda defined}
છે એમ બતાવવા માટે આરંભિક વિધેયો !!{lambda definable} છે અને
!!{lambda definable} આંશિક વિધેયો સંયોજન, આદિમ પુનરાવર્તન
અને અમર્યાદિત શોધ હેઠળ બંધ છે એમ બતાવવું પૂરતું છે.
\end{proof}

સાબિતીનો બાકીનો ભાગ વધુ વાંચનીય બને તે માટે આપણે વધુ પ્રચલિત
સંકેતન વાપરીશું. દાખલા તરીકે, $Mxyz$ને બદલે $M(x, y, z)$ લખીશું.
આ સંકેતન કંઈક સૂચવે છે, પરંતુ યાદ રાખવું કે પ્રકારવિહીન લૅમ્બડા
કલનના પદો સાથે નિશ્ચિત દલીલ-સંખ્યા જોડાયેલી નથી. તેથી એ જ
પદ~$M$ માટે $M(x,y)$ અને $M(x,y,z,w)$ બંને લખવું અર્થપૂર્ણ છે.
તેમ છતાં, આ સંકેતનથી આપણે~$M$ને ત્રણ ચલોનું વિધેય માનીને
વાપરીએ છીએ તે સ્પષ્ટ થાય છે અને વ્યાખ્યાઓ પાછળનો હેતુ સમજાય છે.
એ જ રીતે, ``$M = \lambd[x][\lambd[y][\lambd[z][\dots]]]$
થી~$M$ વ્યાખ્યાયિત કરો'' એમ કહેવાને બદલે ``$M(x,y,z) = \dots$
થી~$M$ વ્યાખ્યાયિત કરો'' એમ કહીશું.

\end{document}
